\documentclass{article}
\usepackage[utf8]{inputenc}
\usepackage{geometry}
\geometry{a4paper, margin=1in}

\title{For My Daughters: A World of Learning Networks}
\author{Albert Jan van Hoek}
\date{\today}

\begin{document}

\maketitle

Today I am thinking about my daughters. About the world they are growing up in. And about what we are creating for them.

For me, the \textbf{Theory of Long-term Collaboration (TLC)} represents hope — that we can actually build something better.

\section*{The Core Idea of TLC}

At the heart of TLC is a simple idea: when we approach collaboration as the interaction between two learning networks, something changes. A learning network is defined by its learning capabilities — not by its outer shell. Whether that network operates in a female body or a male one, whether that body has a certain skin color or sits in a wheelchair, is irrelevant. It is still a learning network. The same goes for the language it speaks. What does matter is how well it learns — what data it is exposed to, what feedback it receives, how openly it communicates.

\section*{The Impact of Information}

Furthermore, from a collaboration perspective, this also means that what passes between networks matters enormously. False information, withheld facts, and denied realities are all damaging to collective learning. They are not just moral failures — they are structurally harmful.

Once we see ourselves as a society of learning networks, we will also think differently about rules and norms. Lying and denying become not just unethical but structurally harmful. We gain a clearer boundary between what undermines collective learning — misinformation, suppression of data — and what causes harm for other reasons, such as patriarchy or racism. It is a different foundation, and I believe a stronger one, a more universal one.

\section*{Beyond Labels}

Which brings me back to my daughters. They are still culturally labelled as women. TLC feels like an antidote to exactly that kind of thinking. Under TLC, they are not defined by their gender, but by who they are as learning beings. And this works both ways — they too will assess others with that same neutrality. Under TLC, those old labels simply lose their grip.

That idea makes me genuinely happy. Not just for them — but for all of us.

\end{document}
